%% =====================================================================
%%  Compute-Matched Post-Training Arena  --  paper skeleton
%%
%%  THIS FILE CONTAINS NO RESULTS. Every number is XX.X or \TODO{...}.
%%  That is deliberate: the numbers must come out of src/ and runs/.
%%
%%  Compile:
%%      pdflatex main && bibtex main && pdflatex main && pdflatex main
%%
%%  Template: ACL 2026 (needs acl.sty + acl_natbib.bst from the ACL
%%  rolling-review template). If you target NeurIPS instead, swap in
%%  neurips_2026.sty and drop the [review] option.
%% =====================================================================
\documentclass[11pt]{article}

%% --- template + package detection ---------------------------------------
%% Every package below is OPTIONAL and probed together with its own
%% dependencies (e.g. graphicx needs grfext, which a minimal MiKTeX may lack).
%% The skeleton therefore compiles on a bare TeX install so you can check the
%% argument structure before the figure/bib pipeline exists.
\newif\ifacltemplate
\IfFileExists{acl.sty}{%
  \acltemplatetrue
  \usepackage[review]{acl}%
}{%
  \acltemplatefalse
}

\ifacltemplate\else
  \IfFileExists{geometry.sty}{\usepackage[margin=1in]{geometry}}{}
  % Deliberately NOT loading natbib on the fallback path. With the 'plain'
  % bibliography style natbib cannot recover author names, so every \citet
  % emits "Author undefined". The \citet/\citep shims further down are enough,
  % and they cannot break on a bare TeX install.
\fi

\IfFileExists{microtype.sty}{\usepackage{microtype}}{}
\IfFileExists{booktabs.sty}{\usepackage{booktabs}}{}
\IfFileExists{color.sty}{%
  \IfFileExists{xcolor.sty}{\usepackage{xcolor}}{}%
}{}
\IfFileExists{grfext.sty}{%
  \IfFileExists{graphicx.sty}{\usepackage{graphicx}}{}%
}{}
\IfFileExists{amsmath.sty}{\usepackage{amsmath}}{}
\IfFileExists{times.sty}{\usepackage{times}}{}
\IfFileExists{T1enc.def}{\usepackage[T1]{fontenc}}{}

%% --- fallbacks: \providecommand is a no-op when the package loaded -------
\providecommand{\toprule}{\hline}
\providecommand{\midrule}{\hline}
\providecommand{\bottomrule}{\hline}
\providecommand{\textcolor}[1]{}
\providecommand{\includegraphics}[2][]{\fbox{PLACEHOLDER FIGURE}}
\providecommand{\citet}[1]{\cite{#1}}
\providecommand{\citep}[1]{\cite{#1}}
\newcommand{\TODO}[1]{\textcolor{red}{[TODO: #1]}}
\newcommand{\XX}{XX.X}

\title{What Do You Buy with One GPU-Hour?\\
       A Compute-Matched Study of Post-Training for Small Language Models}

\author{Anonymous submission}
\date{}

\begin{document}
\maketitle

%% ---------------------------------------------------------------------
\begin{abstract}
Post-training methods for language models---supervised fine-tuning (SFT),
rejection-sampling self-training, preference optimisation, and inference-time
scaling---are usually compared under their own authors' compute budgets, which
makes the reported gains hard to attribute. We instead fix the budget and ask
a practitioner's question: \emph{at the scale of a single consumer GPU, which
post-training investment returns the most accuracy per GPU-hour?} We build a
controlled arena over two open model sizes (0.5B and 1.5B parameters), five
methods, four budget points from 5 to 60 minutes, and two seeds, and we reuse
a \emph{single} sampling pass across three of the methods so that the sampling
cost is charged exactly once. \TODO{one-sentence headline finding, e.g.\
``We find that self-consistency decoding dominates trained post-training on
mathematical reasoning at every budget below \XX{} GPU-hours, while SFT remains
the only method that reliably improves format compliance.''} We release the
harness, all adapters, and the raw per-run logs.
\end{abstract}

%% ---------------------------------------------------------------------
\section{Introduction}

A practitioner with one consumer GPU and a week of wall-clock time faces a
concrete allocation problem: spend the time training an adapter, generating and
voting over samples at inference, or collecting better data? The post-training
literature is rich in methods but almost always reports each method at the
budget its authors chose, on models far larger than a single 16\,GB card can
hold \cite{rafailov2023dpo,deepseek2025r1,lambert2024tulu3}. That literature
therefore answers ``does method $M$ help?'' but not ``is $M$ the best use of my
next ten hours?''

This paper takes the second question seriously and makes three contributions.

\begin{enumerate}
  \item \textbf{A compute-matched arena.} We compare SFT on gold data,
        rejection-sampling fine-tuning (RFT), DPO, self-consistency decoding,
        and an untrained prompt-only baseline under a \emph{fixed GPU-hour
        budget}, sweeping four budget points. Training stops when the budget is
        exhausted, not when validation accuracy plateaus.
  \item \textbf{A shared-sampling design.} RFT and DPO are both derived from a
        single pool of $K$ samples per prompt, so the sampling cost is paid
        once and charged once. Without this, the arena is three times more
        expensive and stops being reproducible on one GPU.
  \item \textbf{A format-dividend diagnostic.} We report accuracy under three
        answer-extraction parsers of increasing strictness, and show how much
        of each method's apparent gain survives the strictest one.
        \TODO{state the measured dividend.}
\end{enumerate}

We are explicit about what this paper is not: it does not propose a new
algorithm, it does not reach state of the art, and its scale is small. Its
claim is scoped to the regime it measures---small open models, tens of
GPU-hours---which is precisely the regime most practitioners actually occupy.

%% ---------------------------------------------------------------------
\section{Related Work}

\paragraph{Post-training pipelines.}
Instruction tuning and preference optimisation form the standard post-training
stack \cite{ouyang2022instructgpt,rafailov2023dpo,lambert2024tulu3}, with
reinforcement learning on verifiable rewards now central to reasoning models
\cite{deepseek2025r1}; policy-gradient methods follow PPO
\cite{schulman2017ppo}. Some pipelines now insert a \emph{midtraining} stage
between pretraining and post-training \cite{midtraining2026}. Surveys of the
area \cite{posttraining2025survey,deepdive2025reasoning} catalogue the design
space but do not attempt cross-method budget normalisation.

\paragraph{Does RL do more than SFT?}
A contested line of work asks whether RL post-training teaches new capabilities
or merely redistributes what the base model already has
\cite{sftmemorizes2025,rlreasoning2025}. Our RFT and DPO arms are constructed
from the \emph{same} sample pool, which lets the reader separate ``new data''
from ``new objective'' effects that these studies must otherwise confound.

\paragraph{Self-training and rejection sampling.}
STaR \cite{zelikman2022star} and RAFT \cite{dong2023raft} fine-tune on
self-generated correct trajectories; our RFT is the sampling-based variant of
this family, and it shares its sample pool with DPO by construction.

\paragraph{Inference-time scaling.}
Self-consistency \cite{wang2022selfconsistency} and compute-optimal test-time
scaling \cite{snell2024scaling} show that decoding-time compute can substitute
for model scale. We ask the complementary question of whether it substitutes
for \emph{post-training} compute, at matched cost.

\paragraph{Scaling and empirical studies.}
\citet{tan2026scaling} characterise RL post-training scaling on Qwen2.5 from
0.5B to 72B and find a predictive power law plus saturation of RL efficiency
with scale; our work is a budget-normalised complement at the small end.
\citet{quagmires2025} warn that high SFT scores can mislead about downstream RL
behaviour, and \citet{interleaving2025} study SFT/RL interleaving. Neither
normalises by GPU-hours across inference-side alternatives.

\paragraph{Evaluation fragility.}
\citet{safetyrepro2026} show alignment leaderboard ranks are
configuration-conditional. Our three-parser reporting is a small, deliberate
concession to the same concern.

%% ---------------------------------------------------------------------
\section{Experimental Setup}

\subsection{Models and data}
We use Qwen2.5-Instruct at 0.5B and 1.5B \cite{qwen2024qwen25}, in bf16 with
LoRA \cite{hu2021lora} ($r{=}16$, $\alpha{=}32$, all linear projections).
Training prompts are drawn from GSM8K \cite{cobbe2021gsm8k} and the algebra
split of MATH \cite{hendrycks2021math}; evaluation uses the GSM8K test set and
MATH-500. We unify the target format to ``reasoning text, then
\texttt{The answer is \textbackslash boxed\{X\}.}'' so that format compliance is
measurable independently of correctness. We strip GSM8K's inline calculator
annotations (\texttt{<{}<...>{}>}) from the reasoning text, because no tool
evaluates them at inference time and MATH solutions contain none. A 5\% dev
slice of the training pool is carved out before any tuning; the test sets are
never used for any decision.

\subsection{Methods}
\paragraph{Prompt-only.} No training; the instruct model is queried directly.

\paragraph{SFT.} LoRA fine-tuning on gold solutions.

\paragraph{RFT.} We sample $K{=}8$ responses per prompt at
$T{=}0.7$, keep the correct ones (verified against the gold answer), and run
the same SFT recipe on them.

\paragraph{DPO.} From the same pool we form (chosen, rejected) pairs wherever a
prompt yielded both a correct and an incorrect response, and optimise the DPO
objective \cite{rafailov2023dpo} with $\beta{=}0.1$. The reference policy is
the base model with the adapter disabled, so no second weight copy is needed.

\paragraph{Self-consistency (SC).} No training; at inference we sample $K$
responses and take the majority answer \cite{wang2022selfconsistency}.

\subsection{Budget accounting}
\label{sec:budget}
We define
\begin{equation}
  \text{GPU-hours} \;=\; \frac{t_{\text{wall}} \cdot n_{\text{GPU}}}{3600},
\end{equation}
measured around synchronised device work. Sampling and training are
\emph{charged}; evaluation is not, because it is identical across methods and
independent of the method's cost. Data preparation is not charged. This
convention is what makes SC non-free: its entire cost is sampling, and charging
it is the point of the comparison. For each budget point we calibrate
throughput over a short warm-up and then convert the budget into a fixed number
of optimiser steps, so every method at budget $b$ consumes the same $b$
GPU-hours rather than converging independently.

\subsection{Evaluation}
We report accuracy under three parsers: \emph{lenient} (last number in the
output), \emph{standard} (\texttt{\textbackslash boxed\{\}} if present, else
the last number; our main metric), and \emph{strict}
(\texttt{\textbackslash boxed\{\}} only, and only if a non-trivial reasoning
trace precedes it). We additionally report format compliance (the fraction of
outputs from which the parser could extract anything) and mean output length.
Each configuration is run with two seeds; we report per-seed values alongside
means.

%% ---------------------------------------------------------------------
\section{Results}

\subsection{Accuracy as a function of budget}
Table~\ref{tab:main} reports accuracy at each budget point, and
Figure~\ref{fig:curve} plots accuracy against cumulative GPU-hours.
\TODO{insert the number of the crossover point, if any.}

\begin{table}[t]
\centering
\small
\begin{tabular}{lcccc}
\toprule
Method & 5\,m & 15\,m & 30\,m & 60\,m \\
\midrule
Prompt-only & \XX & \XX & \XX & \XX \\
Self-consistency & \XX & \XX & \XX & \XX \\
SFT & \XX & \XX & \XX & \XX \\
RFT & \XX & \XX & \XX & \XX \\
DPO & \XX & \XX & \XX & \XX \\
\bottomrule
\end{tabular}
\caption{GSM8K test accuracy (standard parser, mean of 2 seeds) at fixed
GPU-minute budgets, 1.5B. Generated by \texttt{src/make\_tables.py}.}
\label{tab:main}
\end{table}

\begin{figure}[t]
\centering
\IfFileExists{figs/curve_1p5b.pdf}{%
  \includegraphics[width=\linewidth]{figs/curve_1p5b.pdf}%
}{%
  \fbox{\parbox[c][4cm][c]{0.85\linewidth}{\centering
    \textbf{PLACEHOLDER}\\ Generate \texttt{figs/curve\_1p5b.pdf} from
    \texttt{runs/evals/*.json}\\ (x: GPU-hours, y: accuracy, band: seed range)}}%
}
\caption{Accuracy vs.\ cumulative GPU-hours, 1.5B. Shaded band: range across
seeds. \TODO{replace with the real figure; regenerate from
\texttt{runs/evals/*.json}.}}
\label{fig:curve}
\end{figure}

\subsection{The format dividend}
\TODO{Table 3: lenient vs.\ standard vs.\ strict accuracy per method. State how
much of each method's gain is extractable-answer rate rather than correctness.}

\subsection{Where the budget is best spent}
\TODO{Interpret the crossover. If SC dominates below some budget $b^*$ and
training-side methods dominate above it, report $b^*$ with a confidence
interval over seeds. If no crossover exists, say so plainly and report the
constant gap.}

%% ---------------------------------------------------------------------
\section{Threats to Validity}
\paragraph{Single scale above 0.5B.} Our two model sizes bracket a narrow range.
We do not claim the ordering holds at 7B; \citet{tan2026scaling} find RL
efficiency saturates with scale, which if anything favours the inference-side
baselines at larger sizes, but we have not measured it.

\paragraph{Two seeds.} Seed variance is comparable to some of the gaps we
report. We therefore report per-seed numbers rather than only means, and we
avoid significance claims that two seeds cannot support.

\paragraph{One task family.} Mathematical reasoning has a verifiable reward,
which favours SC and RFT and disadvantages DPO. We include GSM8K and MATH-500
but not instruction-following, where we expect the ordering to reverse; that
is a limitation, not a finding.

\paragraph{Hyperparameters.} Matching compute does not match tuning effort. We
give every method the same number of hyperparameter trials (reported in the
appendix), but a method with a better default inherits an advantage we cannot
remove.

%% ---------------------------------------------------------------------
\section{Conclusion}
\TODO{two sentences: what a practitioner with one GPU should do, and what the
result does \emph{not} license them to conclude.}

%% ---------------------------------------------------------------------
%% 'plain' ships with every TeX install; 'plainnat' often does not.
\ifacltemplate\else\bibliographystyle{plain}\fi
\bibliography{refs}

\appendix
\section{Hyperparameters and full run list}
\TODO{auto-generate from \texttt{runs/*/meta.json}: every run, its seed,
budget, step count, learning rate, wall-clock, and measured GPU-hours. Including
the failed runs is what makes Section~\ref{sec:budget} credible.}

\end{document}
